Na matemática, conjuntos são coleções imutáveis de objetos distintos.
Na ciência da computação, conferimos a nomenclatura de \textbf{conjuntos} a qualquer
coleção finita abstrata de objetos.
A imutabilidade, nesse caso, não é um requisito; conjuntos computacionais podem crescer
ou diminuir de tamanho ao longo das operações de um algoritmo que o manipula.
Tais conjuntos que podem ter seu tamanho alterado chamamos de \textbf{dinâmicos}.
Os conjuntos dinâmicos que também permitem o teste de pertinência de um membro dado
são chamados de \textbf{dicionários}.

Os elementos de um conjunto de dados são, geralmente, um registro, cujos campos podem
ser de tipos elementares, ponteiros, ou compostos de outros registros.
Os elementos num conjunto não necessariamente precisam ser todos distintos entre si,
a exemplo da coleção de idades de um grupo de pessoas; nada impede que haja duas com a
mesma idade.
Em alguns casos, o membros de um conjunto possuem uma \textbf{chave} que os identifica, e,
caso as chaves sejam todas distintas, o conjunto pode ser considerado como uma coleção de
chaves.

As operações de um conjunto podem ser de dois tipos: consulta ou modificação.
A primeira apenas resgata um valor, membro ou estatística do conjutno.
A segundo altera o estado do conjunto, como sua ordenação, inserindo, removendo ou
modificando membros, entre outras.
Em geral, quase todo conjunto aceita as seguintes operações:

\paragraph{Insere($S,x$)} Operação que toma um conjunto $S$ e insere o membro apontado
por $x$.
Presume-se que tanto $S$ quanto $x$ já foram inicializados.
Dependendo da aplicação, a operação pode falhar caso um membro $x$ com a mesma chave
já exista em $S$.

\paragraph{Remove($S, x$)} Remove do conjunto $S$ o membro apontado por $x$.

\paragraph{Busca($S, k$)} Retorna um ponteiro para um membro $x$ em $S$ cuja chave seja
igual a $k$, ou $\nil$ caso nenhum elemento contenha tal chave.

Outras operações básicas que acompanham conjuntos dinâmicos são

\paragraph{Máximo($S$)} Retorna o membro $x$ de $S$ cujo valor da chave é máximo no
conjunto

\paragraph{Mínimo($S$)} Retorna o membro $x$ de $S$ cujo valor da chave é mínimo no
conjunto

\paragraph{Predecesor($S, x$)} Num conjunto ordenado, retorna o membro $w \in S$ que
precede $x$, ou $\nil$ quando $x$ é o elemento mínimo.

\paragraph{Sucessor($S, x$)} Num conjunto ordenado, retorna o membro $y \in S$ que
sucede $x$, ou $\nil$ quando $x$ é o elemento máximo.
